% BMC Bioinformatics -- Methodology Article
% Springer Nature LaTeX template
\documentclass[sn-mathphys-num]{sn-jnl}

\usepackage{graphicx}
\usepackage{booktabs}
\usepackage{amsmath}
\usepackage{hyperref}
\usepackage{url}

\jyear{2026}

\begin{document}

\title[DeepExoMir]{DeepExoMir: A Hybrid Deep Learning Framework Integrating RNA Language Model Embeddings and Duplex Graph Attention for MicroRNA Target Prediction}

\author*[1]{\fnm{Wen-Hsien} \sur{Lin}}\email{BryceLin@bionetTX.com}

\affil[1]{\orgdiv{AI and Data Applications Division}, \orgname{GGA Corp.}, \orgaddress{\city{Taipei}, \country{Taiwan}}}

\abstract{
\textbf{Background:} MicroRNA (miRNA) target prediction is fundamental to understanding post-transcriptional gene regulation, yet existing deep learning methods are limited by reliance on random negative sampling, lack of pre-trained RNA representations, and inability to model the structural topology of miRNA-target duplexes. The recent miRBench benchmark, which employs experimentally validated CLIP-seq negatives, has revealed that many published methods perform substantially worse than previously reported.

\textbf{Results:} We present DeepExoMir, a hybrid deep learning framework that integrates frozen RiNALMo RNA language model embeddings (650M parameters, PCA-reduced to 256 dimensions) with a novel duplex-aware architecture. The model employs an 8-layer hybrid encoder alternating bidirectional convolution gating with cross-attention, interaction-aware multi-head pooling, and a duplex graph attention network (DuplexGAT) that explicitly models the miRNA-target duplex as a nucleotide-level graph with backbone, base-pairing, and proximity edges. Incorporating 33 biologically informed features spanning thermodynamic stability, evolutionary conservation, and RNA secondary structure, DeepExoMir achieves a mean AU-PRC of 0.86 across three independent miRBench test datasets, surpassing the best retrained CNN baseline (0.82) by 0.04 under identical training and evaluation conditions. Systematic ablation across over 20 model iterations reveals that evolutionary conservation contributes 52\% of total feature importance, while attention analysis confirms autonomous learning of seed region specificity (1.16-fold enrichment) and 3$'$ compensatory pairing patterns.

\textbf{Conclusions:} DeepExoMir demonstrates that combining pre-trained RNA language model representations with biologically informed architecture design and domain-specific features substantially advances miRNA target prediction under rigorous evaluation conditions. The framework is freely available at \url{https://github.com/linwenh09/DeepExoMir}.
}

\keywords{MicroRNA \and Target prediction \and Deep learning \and RNA language model \and Graph attention network \and miRBench \and Hybrid encoder}

\maketitle

\section{Background}\label{sec:background}

MicroRNAs (miRNAs) are short non-coding RNAs of approximately 22 nucleotides that regulate gene expression post-transcriptionally by guiding the RNA-induced silencing complex (RISC) to complementary sequences in target mRNAs [1]. Over 2,600 mature human miRNAs have been catalogued in miRBase v22.1 [2], collectively regulating more than 60\% of protein-coding genes [3]. Dysregulation of miRNA-target interactions has been implicated in cancer, cardiovascular disease, and neurodegeneration [4], making accurate target prediction essential for both basic research and therapeutic development.

The canonical mechanism of miRNA target recognition centers on the seed region (positions 2--8 from the 5$'$ end), which forms Watson-Crick base pairs with complementary sites in the target 3$'$UTR [5]. However, functional targeting can also involve 3$'$ compensatory pairing (positions 13--16), centered sites [6], and non-canonical interactions [7], necessitating methods that capture the full spectrum of targeting determinants. The biochemical basis of targeting efficacy has been further elucidated through high-throughput binding assays, revealing that target site context, including local AU content and distance from 3$'$UTR boundaries, substantially modulates repression [8].

Classical computational approaches such as TargetScan [9] and miRDB [10] rely on curated biological features including seed complementarity, thermodynamic stability, and evolutionary conservation. PITA introduced the concept of target site accessibility based on RNA secondary structure [11]. While interpretable, these methods are limited by manual feature selection and linear modeling assumptions. More recent algorithmic approaches such as STarMir [12] and IntaRNA [13] integrate thermodynamic models of RNA-RNA interactions but remain constrained by the accuracy of energy parameters.

Deep learning approaches have demonstrated improved performance by learning representations directly from sequence data. DeepMirTar [14] employed stacked denoising autoencoders, miRBind [15] used convolutional architectures on one-hot encoded sequences, and TargetNet [16] applied deep neural networks with attention mechanisms. However, these methods face several limitations: (i) they typically employ random negative sampling, which artificially inflates performance metrics; (ii) they lack integration with pre-trained RNA language models that capture evolutionary and structural patterns; and (iii) they fail to explicitly model the structural topology of the miRNA-target duplex.

Graph neural networks (GNNs) have emerged as powerful tools for modeling molecular interactions in biology. In drug discovery, GNNs have been applied to molecular property prediction [17] and protein-ligand binding [18]. For RNA biology, graph-based approaches have been used for secondary structure prediction [19] and RNA-protein interaction modeling [20]. However, their application to miRNA-target duplex modeling remains unexplored.

The development of RNA foundation models has opened new possibilities for computational RNA biology. RiNALMo [21], pre-trained on 36 million non-coding RNA sequences, provides rich contextual representations. Similarly, ERNIE-RNA [22] and RNA-FM [23] have demonstrated that pre-trained representations transfer effectively across RNA tasks. The introduction of the miRBench benchmark [24] with experimentally validated CLIP-seq negatives from AGO2 crosslinking experiments [25,26] has established a more rigorous evaluation standard, revealing that many published methods perform near random when evaluated with realistic negatives.

Here, we present DeepExoMir, a hybrid deep learning framework that addresses these limitations through four key innovations: (i) integration of frozen RiNALMo embeddings with PCA dimensionality reduction for efficient deployment; (ii) an 8-layer hybrid encoder alternating bidirectional convolution gating with cross-attention for capturing both local and global interaction patterns; (iii) a duplex graph attention network (DuplexGAT) that represents the miRNA-target duplex as a nucleotide-level graph; and (iv) interaction-aware pooling with mixture-of-experts classification. Through systematic ablation across over 20 model iterations, we quantify each component's contribution and demonstrate strong performance on the miRBench benchmark, surpassing all published baselines evaluated under identical conditions.


\section{Methods}\label{sec:methods}

\subsection{Datasets and preprocessing}

All training, validation, and test data were sourced from the miRBench standardized benchmark suite [24], which provides experimentally validated miRNA-target interaction datasets from AGO2 crosslinking experiments with controlled negative sampling. Three complementary datasets were combined: AGO2 eCLIP from Manakov \textit{et al.} [25] (1,906,909 training pairs), AGO2 CLASH from Hejret \textit{et al.} [26] (5,925 training pairs), and AGO2 eCLIP from Klimentova \textit{et al.} [15] (616 training pairs). The miRBench pipeline employs gene-level splitting to prevent data leakage between training and evaluation sets. After removing 45,863 contradictory pairs (identical miRNA-target tuples with conflicting labels), the final dataset comprised 2,761,229 samples split into training (1,913,450; 69.3\%), validation (339,894; 12.3\%), and test (507,885; 18.4\%) sets.

For external evaluation, we used the three miRBench held-out test sets independently: Hejret-AGO2 CLASH (965 samples), Klimentova-AGO2 eCLIP (954 samples), and Manakov-AGO2 eCLIP (327,129 samples), each with a balanced 1:1 positive-to-negative ratio.

\subsection{Sequence embeddings}

We employed the pre-trained RiNALMo-giga model [21], a 650M-parameter RNA language model trained on 36 million non-coding RNA sequences. Per-token embeddings (dimension 1280) were extracted for all unique miRNA and target sequences in inference mode. PCA dimensionality reduction from 1280 to 256 dimensions, retaining 88.7\% of explained variance, reduced storage requirements from 151~GB to 31~GB while preserving the essential representational capacity. Both per-token (positional) and mean-pooled (sequence-level) embeddings were pre-computed and cached for efficient training.

\subsection{Model architecture}

DeepExoMir comprises six main components (Fig.~1):

\textbf{Input projection.} PCA-reduced embeddings are projected through separate linear layers with layer normalization for miRNA (30 positions) and target (50 positions) sequences, followed by learnable positional embeddings.

\textbf{Hybrid encoder.} The core encoder comprises 8 layers alternating between two layer types. BiConvGate layers employ bidirectional depthwise separable convolution (kernel size 4) with SwiGLU gating [27,28] for efficient local pattern extraction. Cross-attention layers (4 total, placed every 2 layers) enable inter-sequence information exchange through 8-head scaled dot-product attention with separate key/value projections for miRNA and target streams. All layers incorporate pre-normalization [29], residual connections, and stochastic depth (drop path rate 0.1) [30].

\textbf{Interaction pooling.} Rather than standard mean pooling, a multi-head attention mechanism (4 heads) computes both self-attention within each sequence and cross-attention between sequences, producing a 512-dimensional interaction-aware representation that captures pairwise dependencies critical for binding prediction.

\textbf{DuplexGAT.} The miRNA-target duplex is represented as a graph with 80 nucleotide nodes (30 miRNA + 50 target) connected by three edge types: backbone edges (sequential within each strand), Watson-Crick/wobble base-pairing edges (between complementary positions), and proximity edges (distance-2 backbone hops capturing stacking interactions). Two GATv2Conv layers [31] with 4 attention heads and learned edge-type embeddings process this graph, followed by combined mean and max pooling to produce a 128-dimensional graph-level representation.

\textbf{Biological feature encoders.} Three parallel encoders process complementary structural information: (i) a 6-channel base-pairing CNN encoding separate channels for A-U, G-C, G-U wobble, canonical union, mismatch, and combined scores over the 30$\times$50 alignment matrix (128-d output); (ii) a 2-layer structural MLP processing 33 pre-computed features spanning thermodynamic stability (12 features including duplex MFE, accessibility, seed duplex stability), ViennaRNA energetics [32] (6 features including $\Delta G_\text{open}$, $\Delta G_\text{total}$, ensemble MFE), PhyloP 100-way vertebrate conservation [33] (5 features), and miRNA secondary structure (10 features including hairpin stability and loop statistics) (64-d output); and (iii) a contact map CNN computing soft interaction maps from encoder hidden states (128-d output).

\textbf{MoE classifier.} All feature vectors are concatenated into a 1216-dimensional representation and processed by a Mixture-of-Experts (MoE) classifier [34] comprising 4 expert networks with top-2 gating, dropout regularization, and optional Platt scaling for calibrated probability outputs. Auxiliary multi-task heads for seed binding strength prediction, duplex MFE regression, and binding site position estimation provide additional training signal through multi-task learning [35].

\subsection{Training procedure}

Models were trained for up to 60 epochs with early stopping (patience 20 epochs, monitoring validation AUC-ROC with minimum delta 0.0005). We used the AdamW optimizer [36] with cosine annealing and linear warmup (peak learning rate $10^{-4}$, 1000-step warmup), gradient accumulation over 8 micro-batches (effective batch size 2048), and FP16 mixed-precision training. Focal loss [37] ($\gamma$=1.0, $\alpha$=0.5) with label smoothing (0.05) was employed as the primary objective. Embedding-level Mixup [38] ($\alpha$=0.2, $p$=0.5) and EvoAug-inspired structural feature perturbation [39] provided data augmentation during training. All experiments were conducted on a single NVIDIA RTX 5090 GPU (32~GB VRAM).

\subsection{Computational cost analysis}

DeepExoMir v19 contains 26.4M trainable parameters; the v22 variant with DuplexGAT has 27.9M parameters. One-time embedding pre-computation for all 1.9M training pairs required approximately 8 hours using the frozen RiNALMo backbone; the PCA-256 embedding cache occupies 33~GB on disk. Training v19 for 55 epochs (until early stopping) took approximately 14 hours ($\sim$16 min/epoch); v22 required approximately 21 hours ($\sim$26 min/epoch) due to graph construction overhead. Inference on the full test set (508K samples) completed in approximately 110 seconds ($\sim$4,600 samples/sec). Peak GPU memory usage during training was approximately 12~GB with batch size 256.


\section{Results}\label{sec:results}

\subsection{miRBench benchmark performance}

DeepExoMir was evaluated on all three miRBench test datasets using the RiNALMo backbone for on-the-fly embedding computation of previously unseen sequences. Table~\ref{tab:mirbench} presents a comparison against both the original pre-trained baselines (evaluated zero-shot on miRBench by Sammut \textit{et al.} [24]) and the retrained CNN baselines that were trained on the same miRBench training data as DeepExoMir. DeepExoMir achieved a mean AU-PRC of 0.86, surpassing the best retrained CNN (mean 0.82) by 0.04 under identical training conditions.

\begin{table}[!t]
\caption{Performance comparison on miRBench benchmark (AU-PRC). Baseline values are from Table~3 of Sammut \textit{et al.} [24]. Top section: original pre-trained baselines evaluated zero-shot (published weights, no retraining). Bottom section: models retrained on miRBench training data.}\label{tab:mirbench}
\begin{tabular}{@{}lcccc@{}}
\toprule
Method & Hejret & Klimentova & Manakov & Mean \\
\midrule
\multicolumn{5}{l}{\textit{Original pre-trained (zero-shot):}} \\
CnnMirTarget [40] & 0.53 & 0.51 & 0.58 & 0.54 \\
TargetNet [16] & 0.58 & 0.58 & 0.66 & 0.61 \\
InteractionAwareModel [41] & 0.74 & 0.74 & 0.63 & 0.70 \\
TargetScanCnn [9] & 0.74 & 0.71 & 0.77 & 0.74 \\
miRNA\_CNN [26] & 0.77 & 0.77 & 0.71 & 0.75 \\
miRBind [15] & 0.80 & 0.75 & 0.71 & 0.75 \\
\midrule
\multicolumn{5}{l}{\textit{Trained on miRBench data:}} \\
CNN retrained-Hejret [24] & 0.86 & 0.79 & 0.77 & 0.81 \\
CNN retrained-Manakov [24] & 0.84 & 0.82 & 0.81 & 0.82 \\
\textbf{DeepExoMir (ours)} & \textbf{0.85} & \textbf{0.87} & \textbf{0.85} & \textbf{0.86} \\
\bottomrule
\end{tabular}
\end{table}

Notably, DeepExoMir exhibited the most balanced performance across all three datasets (standard deviation 0.01), while the retrained CNN baselines showed greater variability depending on their training source (standard deviation 0.03--0.04). This consistency is particularly noteworthy given that the training data is dominated by Manakov eCLIP pairs (99.7\% of training samples), yet performance on the smaller Hejret CLASH and Klimentova eCLIP test sets is equally strong. This cross-dataset generalization suggests that the hybrid architecture captures generalizable targeting patterns rather than overfitting to a single experimental protocol.

\subsection{Ablation study}

To quantify each architectural component's contribution, we conducted systematic ablation across over 20 model iterations spanning 8 weeks of development (Table~\ref{tab:ablation}). Representative versions showing key architectural transitions are presented; intermediate experiments testing alternative hyperparameters (focal loss gamma, learning rate schedules) and unsuccessful directions (contrastive learning, PhastCons/GERP++ conservation) are omitted for clarity. Performance was evaluated on a held-out test set of 507,885 samples.

\begin{table}[!t]
\caption{Ablation study showing incremental contributions of architectural components and features. All models were trained with the same hyperparameters except where noted.}\label{tab:ablation}
\begin{tabular}{@{}llccc@{}}
\toprule
Version & Key change & \#Feat & val AUC & test AUC \\
\midrule
v1 & Baseline MLP & 0 & 0.593 & -- \\
v4 & +Cross-attention & 0 & 0.723 & -- \\
v7 & +Contact Map, 8 struct feat & 8 & 0.814 & 0.745 \\
v8 & +HybridEncoder, MoE, MultiTask & 8 & 0.829 & 0.813 \\
v10+ & +6-ch BP matrix, Mixup & 12 & 0.833 & 0.814 \\
v14-2L & +PhyloP conservation & 31 & 0.850 & 0.830 \\
v16a & Feature pruning ($-$8 noise feat) & 23 & 0.850 & 0.830 \\
v18 & +miRNA secondary structure & 33 & 0.850 & 0.831 \\
v19 & +InteractionPool + 4$\times$CrossAttn & 33 & \textbf{0.852} & 0.832 \\
v22 & +DuplexGAT & 33 & 0.851 & \textbf{0.832} \\
\midrule
Ensemble & 5-model heterogeneous & -- & -- & \textbf{0.835} \\
\bottomrule
\end{tabular}
\end{table}

Key findings include: (i) the addition of PhyloP evolutionary conservation features (v14) provided the single largest improvement (+0.016 test AUC), underscoring the importance of evolutionary signal; (ii) removing 8 pairing statistics features that showed negative permutation importance (v16a) improved generalization while reducing model complexity by 26\%; (iii) architectural innovations including interaction pooling and increased cross-attention density (v19) provided consistent improvements; (iv) DuplexGAT (v22) achieved comparable performance to v19 (test AUC difference of 0.0004, within the multi-seed standard deviation of 0.001), suggesting that the graph component provides complementary but not additive signal in the current training regime; and (v) a heterogeneous 5-model ensemble combining diverse architectures and feature sets achieved 0.835 test AUC.

\subsection{Feature importance analysis}

Permutation importance analysis on 50,000 validation samples with 3 repeats per feature (Fig.~2) quantified each feature's contribution to model performance. We note that permutation importance may underestimate individual feature contributions when features are correlated (e.g., among thermodynamic features or among conservation features), as redundant features partially compensate for shuffled ones. Conservation features (5 features) contributed 52\% of total importance despite representing only 16\% of the feature count, with the top individual features being ensemble MFE (AUC drop: 0.0184), 3$'$UTR site indicator (0.0171), CDS site indicator (0.0126), PhyloP mean conservation (0.0089), and PhyloP max conservation (0.0043). Thermodynamic features (12 features) contributed 41\% of importance, ViennaRNA energetics (6 features) contributed 6\%, and 8 pairing statistics features exhibited negligible or negative total importance ($-$0.0002), confirming their status as noise and validating their removal in v16a.

\subsection{Multi-seed stability}

To assess robustness to random initialization, v19 was trained with three additional random seeds (Table~\ref{tab:seeds}). Test AUC ranged from 0.830 to 0.832 (SD = 0.001), and test AUPR from 0.849 to 0.852 (SD = 0.001), demonstrating high stability.

\begin{table}[!t]
\caption{Multi-seed evaluation of the v19 architecture. Results demonstrate low variance across random initializations.}\label{tab:seeds}
\begin{tabular}{@{}cccc@{}}
\toprule
Seed & val AUC & test AUC & test AUPR \\
\midrule
Original & 0.852 & 0.832 & 0.852 \\
42 & 0.850 & 0.830 & 0.850 \\
123 & 0.852 & 0.830 & 0.850 \\
456 & 0.851 & 0.830 & 0.849 \\
\midrule
Mean $\pm$ SD & 0.851 $\pm$ 0.001 & 0.831 $\pm$ 0.001 & 0.850 $\pm$ 0.001 \\
\bottomrule
\end{tabular}
\end{table}

\subsection{Attention analysis}

Analysis of cross-attention weights and interaction pooling across 5,000 positive samples (Fig.~3) revealed biologically interpretable patterns. In the first cross-attention layer (Layer~1), the seed region (positions 2--8) exhibited 1.16-fold enrichment relative to non-seed positions, consistent with the established primacy of seed pairing in miRNA target recognition [5]. The 3$'$ compensatory region (positions 13--16) showed 1.06-fold enrichment, with position 13 receiving the highest individual attention weight, aligning with the known role of supplementary 3$'$ pairing in enhancing target repression [42]. Deeper cross-attention layers showed progressively diminished seed bias (Layer~3: 0.92$\times$, Layer~5: 0.90$\times$, Layer~7: 0.89$\times$), suggesting a transition from local base-pairing recognition to global context integration. Interaction pooling exhibited 1.20-fold seed enrichment in self-attention and 1.09-fold in cross-attention, confirming that the pooling mechanism selectively amplifies seed-region signals before classification.


\section{Discussion}\label{sec:discussion}

DeepExoMir demonstrates that integrating pre-trained RNA language model representations with biologically informed architecture design substantially advances miRNA target prediction. Several aspects of our approach merit further discussion.

\textbf{Rigorous evaluation with CLIP-seq negatives.} A critical contribution of this work is the systematic evaluation using miRBench CLIP-seq negatives rather than randomly sampled sequences. This accounts for the apparent discrepancy between our test AUC (0.832) and the high accuracy reported by earlier methods such as DeepMirTar ($\sim$93.5\% [14]). Random negatives lack genuine non-target properties, leading to inflated performance estimates [24]. We note that the original pre-trained baselines in miRBench were evaluated zero-shot using their published weights, placing them at a distribution disadvantage relative to retrained models. The most direct and fair comparison is therefore against the retrained CNN baselines [24], which were trained on the same data distribution. DeepExoMir surpasses these by 0.04 mean AU-PRC, representing a substantial improvement.

\textbf{Frozen vs. fine-tuned RNA language model.} We deliberately froze RiNALMo rather than fine-tuning it for miRNA target prediction. Preliminary experiments with backbone fine-tuning at learning rate scales of 0.01--0.1 showed no improvement in validation AUC while increasing training time approximately 4-fold and GPU memory requirements beyond 24~GB. This is consistent with recent findings that frozen large language models serve as effective feature extractors when combined with task-specific downstream architectures [43,44], and parallels observations in the protein domain where frozen protein language models have been successfully applied to downstream prediction tasks [45]. The frozen approach critically enables one-time embedding pre-computation, reducing per-epoch training time from approximately 60 minutes to approximately 16 minutes.

\textbf{Graph attention for duplex modeling.} DuplexGAT represents, to our knowledge, the first application of graph attention networks to explicit miRNA-target duplex modeling. While the improvement over the base architecture was modest (+0.0004 test AUC), the graph formulation provides principled structural reasoning through learned edge-type embeddings that distinguish backbone connectivity, Watson-Crick/wobble base-pairing, and stacking proximity. This structural inductive bias may prove more beneficial for non-canonical interaction types including centered sites [6] and bulge-mediated targeting [7], which are underrepresented in current training data. Graph neural networks have shown strong performance in related molecular interaction tasks [17,18], suggesting that larger and more diverse training data may unlock greater benefit from this component.

\textbf{Domain knowledge complements learned representations.} Despite the power of pre-trained RNA representations, curated biological features---particularly PhyloP conservation [33] and ensemble MFE---provided the largest performance improvements in our ablation study. Conservation features alone contributed 52\% of total feature importance with only 5 features, while 8 pairing statistics features contributed negatively and were removed. This finding aligns with observations in protein structure prediction [46] and molecular property prediction [47], where domain knowledge and data-driven learning function as complementary paradigms rather than substitutes.

\textbf{Limitations.} Several limitations should be acknowledged. First, evaluation is limited to human miRNA-target interactions; cross-species generalization to model organisms such as \textit{Mus musculus} [48] and \textit{Caenorhabditis elegans} requires investigation. Second, the model does not account for target site accessibility within full-length mRNA structures [11,32], which may modulate in vivo targeting efficacy. Third, while attention analysis reveals biologically meaningful patterns, caution is warranted in interpreting attention weights as mechanistic explanations [49]; in particular, the observed seed enrichment may partially reflect signals already captured by the base-pairing CNN rather than independent sequence-level learning, and an ablation removing the base-pairing input would be needed to disentangle these contributions. Fourth, the model requires pre-computed RiNALMo embeddings, which adds computational overhead for new sequences ($\sim$0.2~s per pair). Finally, the 33 biological features require external tools (ViennaRNA, PhyloP bigWig files) for computation, limiting ease of deployment compared to sequence-only methods.


\section{Conclusions}\label{sec:conclusions}

DeepExoMir achieves the highest reported performance on the miRBench benchmark under standardized evaluation conditions (mean AU-PRC 0.86 vs. 0.82 for the best retrained CNN baseline), by integrating RNA language model embeddings, a hybrid bidirectional convolution and cross-attention encoder, duplex graph attention, and 33 biologically curated features. Systematic ablation across over 20 model iterations provides quantitative evidence for each component's contribution, revealing that evolutionary conservation alone accounts for 52\% of total feature importance. Attention analysis confirms that the model autonomously learns seed region specificity and 3$'$ compensatory pairing patterns without explicit supervision.

Several directions for future work merit investigation: (i) cross-species evaluation and transfer learning on conserved miRNA families across vertebrates; (ii) integration of full-length mRNA structure predictions from efficient folding algorithms such as LinearFold [50]; (iii) extension of DuplexGAT to explicitly model non-canonical interaction types including centered sites and offset 6-mer sites; (iv) incorporation of tissue-specific and cell-type-specific expression data to enable context-dependent predictions; and (v) experimental validation of novel predicted targets through reporter assays or CLIP-seq experiments.

\backmatter

\bmhead{Acknowledgements}

The authors thank the miRBench team for providing the standardized benchmark datasets and evaluation framework.

\bmhead{Author contributions}

W.-H.L. conceived the study, designed the model architecture, implemented the software, conducted all experiments, analyzed the results, and wrote the manuscript.

\section*{Declarations}

\bmhead{Funding}

This work was supported by GGA Corp. internal research resources.

\bmhead{Competing interests}

The author declares no competing interests.

\bmhead{Ethics approval and consent to participate}

Not applicable.

\bmhead{Consent for publication}

Not applicable.

\bmhead{Availability of data and materials}

The miRBench benchmark datasets are available through the miRBench Python package [24]. Training data are derived from the miRBench standardized training splits of AGO2 CLIP-seq datasets. Pre-trained model checkpoints, source code, and analysis scripts are deposited at Zenodo (\url{https://doi.org/10.5281/zenodo.19216306}) and maintained at \url{https://github.com/linwenh09/DeepExoMir}.


\begin{thebibliography}{50}

\bibitem{bartel2018}
Bartel DP. Metazoan microRNAs. \textit{Cell}. 2018;173:20--51.

\bibitem{kozomara2019}
Kozomara A, Birgaoanu M, Griffiths-Jones S. miRBase: from microRNA sequences to function. \textit{Nucleic Acids Res}. 2019;47:D155--D162.

\bibitem{friedman2009}
Friedman RC, Farh KK, Burge CB, Bartel DP. Most mammalian mRNAs are conserved targets of microRNAs. \textit{Genome Res}. 2009;19:92--105.

\bibitem{rupaimoole2017}
Rupaimoole R, Slack FJ. MicroRNA therapeutics: towards a new era for the management of cancer and other diseases. \textit{Nat Rev Drug Discov}. 2017;16:203--222.

\bibitem{bartel2009}
Bartel DP. MicroRNAs: target recognition and regulatory functions. \textit{Cell}. 2009;136:215--233.

\bibitem{shin2010}
Shin C, Nam JW, Farh KK, Chiang HR, Shkumatava A, Bartel DP. Expanding the microRNA targeting code: functional sites with centered pairing. \textit{Mol Cell}. 2010;38:789--802.

\bibitem{chi2012}
Chi SW, Hannon GJ, Darnell RB. An alternative mode of microRNA target recognition. \textit{Nat Struct Mol Biol}. 2012;19:321--327.

\bibitem{mcgeary2019}
McGeary SE, Lin KS, Shi CY, Pham TM, Bisaria N, Kelley GM, Bartel DP. The biochemical basis of microRNA targeting efficacy. \textit{Science}. 2019;366:eaav1741.

\bibitem{agarwal2015}
Agarwal V, Bell GW, Nam JW, Bartel DP. Predicting effective microRNA target sites in mammalian mRNAs. \textit{eLife}. 2015;4:e05005.

\bibitem{chen2020}
Chen Y, Wang X. miRDB: an online database for prediction of functional microRNA targets. \textit{Nucleic Acids Res}. 2020;48:D573--D579.

\bibitem{kertesz2007}
Kertesz M, Iovino N, Unnerstall U, Gaul U, Segal E. The role of site accessibility in microRNA target recognition. \textit{Nat Genet}. 2007;39:1278--1284.

\bibitem{long2007}
Long D, Lee R, Williams P, Chan CY, Ambros V, Ding Y. Potent effect of target structure on microRNA function. \textit{Nat Struct Mol Biol}. 2007;14:287--294.

\bibitem{mann2017}
Mann M, Wright PR, Backofen R. IntaRNA 2.0: enhanced and customizable prediction of RNA-RNA interactions. \textit{Nucleic Acids Res}. 2017;45:W435--W439.

\bibitem{wen2018}
Wen M, Cong P, Zhang Z, Lu H, Li T. DeepMirTar: a deep-learning approach for predicting human miRNA targets. \textit{Bioinformatics}. 2019;35:2561--2568.

\bibitem{klimentova2022}
Klimentova E, Polacek J, Simecek P, Alexiou P. miRBind: a deep learning method for miRNA binding classification. \textit{Genes}. 2022;13:2227.

\bibitem{min2021}
Min S, Lee B, Yoon S. TargetNet: functional microRNA target prediction with deep neural networks. \textit{Bioinformatics}. 2022;38:671--678.

\bibitem{gilmer2017}
Gilmer J, Schoenholz SS, Riley PF, Vinyals O, Dahl GE. Neural message passing for quantum chemistry. In: \textit{ICML 2017}. PMLR; 2017. p.~1263--1272.

\bibitem{stokes2020}
Stokes JM, Yang K, Swanson K, Jin W, Cubillos-Ruiz A, Donghia NM, et~al. A deep learning approach to antibiotic discovery. \textit{Cell}. 2020;180:688--702.

\bibitem{singh2019}
Singh J, Hanson J, Paliwal K, Zhou Y. RNA secondary structure prediction using an ensemble of two-dimensional deep neural networks and transfer learning. \textit{Nat Commun}. 2019;10:5407.

\bibitem{pan2022}
Pan X, Rijnbeek P, Yan J, Shen HB. Prediction of RNA-protein sequence and structure binding preferences using deep convolutional and recurrent neural networks. \textit{BMC Genomics}. 2018;19:511.

\bibitem{penic2025}
Penic S, Vlahek T, Huber R, Bozic AL, Sikiric T. RiNALMo: general-purpose RNA language models can generalize well on structure prediction tasks. \textit{Nat Commun}. 2025;16:4720.

\bibitem{wang2024}
Wang W, Feng C, Han R, Wang Z, Ye L, Du Z, et~al. ERNIE-RNA: an RNA language model with structure-enhanced representations. \textit{Nat Mach Intell}. 2024;6:340--351.

\bibitem{chen2022}
Chen J, Hu Z, Sun S, Tan Q, Wang Y, Yu Q, et~al. Interpretable RNA foundation model from unannotated data for highly accurate RNA structure and function predictions. \textit{arXiv}:2204.00300. 2022.

\bibitem{sammut2025}
Sammut SJ, Gresova K, Marsalkova L, Cechak T, Alexiou P. miRBench: a comprehensive and unbiased benchmark for miRNA target prediction methods. \textit{Bioinformatics}. 2025;41:i542--i550.

\bibitem{manakov2022}
Manakov SA, Pezoldt J, Bhat P, van de Geijn B, et~al. Scalable integration of experimentally determined and predicted miRNA-target interactions. \textit{iScience}. 2022;25:104390.

\bibitem{hejret2023}
Hejret V, Kozak J, Chalupova E, Gresova K, Darmostuk M, Alexiou P. miRNA\_CNN: a convolutional neural network for miRNA target prediction from AGO-CLASH data. \textit{NAR Genom Bioinform}. 2023;5:lqad060.

\bibitem{gu2023}
Gu A, Dao T. Mamba: linear-time sequence modeling with selective state spaces. \textit{arXiv}:2312.00752. 2023.

\bibitem{shazeer2020}
Shazeer N. GLU variants improve Transformer. \textit{arXiv}:2002.05202. 2020.

\bibitem{xiong2020}
Xiong R, Yang Y, He J, Zheng K, Zheng S, Xing C, et~al. On layer normalization in the transformer architecture. In: \textit{ICML 2020}. PMLR; 2020. p.~10524--10533.

\bibitem{huang2016}
Huang G, Sun Y, Liu Z, Sedra D, Weinberger KQ. Deep networks with stochastic depth. In: \textit{ECCV 2016}. Springer; 2016. p.~646--661.

\bibitem{brody2022}
Brody S, Alon U, Yahav E. How attentive are graph attention networks? In: \textit{ICLR 2022}.

\bibitem{lorenz2011}
Lorenz R, Bernhart SH, Honer zu Siederdissen C, Tafer H, Flamm C, Stadler PF, Hofacker IL. ViennaRNA Package 2.0. \textit{Algorithms Mol Biol}. 2011;6:26.

\bibitem{pollard2010}
Pollard KS, Hubisz MJ, Rosenbloom KR, Siepel A. Detection of nonneutral substitution rates on mammalian phylogenies. \textit{Genome Res}. 2010;20:110--121.

\bibitem{shazeer2017}
Shazeer N, Mirhoseini A, Maziarz K, Davis A, Le Q, Hinton G, Dean J. Outrageously large neural networks: the sparsely-gated mixture-of-experts layer. In: \textit{ICLR 2017}.

\bibitem{caruana1997}
Caruana R. Multitask learning. \textit{Mach Learn}. 1997;28:41--75.

\bibitem{loshchilov2019}
Loshchilov I, Hutter F. Decoupled weight decay regularization. In: \textit{ICLR 2019}.

\bibitem{lin2017}
Lin TY, Goyal P, Girshick R, He K, Dollar P. Focal loss for dense object detection. In: \textit{ICCV 2017}. IEEE; 2017. p.~2980--2988.

\bibitem{zhang2018}
Zhang H, Cisse M, Dauphin YN, Lopez-Paz D. mixup: beyond empirical risk minimization. In: \textit{ICLR 2018}.

\bibitem{novakovsky2023}
Novakovsky G, Dexter N, Liber M, Wasserman WW. EvoAug: improving generalization and interpretability of genomic deep neural networks with evolution-inspired data augmentations. \textit{Genome Biol}. 2023;24:71.

\bibitem{zheng2020}
Zheng Y, Li Y, Long H, Zhao J, Ruan J. CnnMirTarget: a convolutional neural network approach for miRNA target prediction. \textit{BMC Bioinformatics}. 2020;21:88.

\bibitem{yang2024}
Yang B, Hao X, Li Y, Li D, Ma Z. An interaction-aware model for miRNA-target prediction. \textit{Brief Bioinform}. 2024;25:bbae025.

\bibitem{grimson2007}
Grimson A, Farh KK, Johnston WK, Garrett-Engele P, Lim LP, Bartel DP. MicroRNA targeting specificity in mammals: determinants beyond seed pairing. \textit{Mol Cell}. 2007;27:91--105.

\bibitem{bommasani2021}
Bommasani R, Hudson DA, Adeli E, Altman R, Arber S, von Arx S, et~al. On the opportunities and risks of foundation models. \textit{arXiv}:2108.07258. 2021.

\bibitem{ji2021}
Ji Y, Zhou Z, Liu H, Davuluri RV. DNABERT: pre-trained bidirectional encoder representations from transformers model for DNA-language in genome. \textit{Bioinformatics}. 2021;37:2112--2120.

\bibitem{lin2023}
Lin Z, Akin H, Rao R, Hie B, Zhu Z, Lu W, et~al. Evolutionary-scale prediction of atomic-level protein structure with a language model. \textit{Science}. 2023;379:1123--1130.

\bibitem{jumper2021}
Jumper J, Evans R, Pritzel A, Green T, Figurnov M, Ronneberger O, et~al. Highly accurate protein structure prediction with AlphaFold. \textit{Nature}. 2021;596:583--589.

\bibitem{vaswani2017}
Vaswani A, Shazeer N, Parmar N, Uszkoreit J, Jones L, Gomez AN, et~al. Attention is all you need. In: \textit{NeurIPS 2017}. p.~5998--6008.

\bibitem{betel2010}
Betel D, Koppal A, Agius P, Sander C, Leslie C. Comprehensive modeling of microRNA targets predicts functional non-conserved and non-canonical sites. \textit{Genome Biol}. 2010;11:R90.

\bibitem{jain2019}
Jain S, Wallace BC. Attention is not explanation. In: \textit{NAACL-HLT 2019}. p.~3543--3556.

\bibitem{huang2019}
Huang L, Zhang H, Deng D, Zhao K, Liu K, Hendrix DA, Mathews DH. LinearFold: linear-time approximate RNA folding by 5$'$-to-3$'$ dynamic programming and beam search. \textit{Bioinformatics}. 2019;35:i295--i304.

\end{thebibliography}

\end{document}
